\documentclass[10pt,a4paper]{amsart}
\usepackage{amsmath,amssymb,amsthm,graphicx}
% Retained minimal preamble for the Stacks Project chapter sites-cohomology.tex (Cohomology on Sites).
% Verbatim excerpt of the chapter preamble: Xy-pic setup (native preamble lines 14-18) and the native
% theorem declarations with the mathematical macros (native preamble lines 154-217). Copied mechanically;
% nothing added, removed or redefined. The wrapper class and the amsmath/amssymb/amsthm/graphicx base
% packages are supplied by the assembler, so the class/externaldocument part of the native preamble is omitted.
% For commutative diagrams we use Xy-pic
\usepackage[all]{xy}

% We use 2cell for 2-commutative diagrams.
\xyoption{2cell}
% Theorem environments.
%
\theoremstyle{plain}
\newtheorem{theorem}[subsection]{Theorem}
\newtheorem{proposition}[subsection]{Proposition}
\newtheorem{lemma}[subsection]{Lemma}

\theoremstyle{definition}
\newtheorem{definition}[subsection]{Definition}
\newtheorem{example}[subsection]{Example}
\newtheorem{exercise}[subsection]{Exercise}
\newtheorem{situation}[subsection]{Situation}

\theoremstyle{remark}
\newtheorem{remark}[subsection]{Remark}
\newtheorem{remarks}[subsection]{Remarks}

\numberwithin{equation}{subsection}

% Macros
%
\def\lim{\mathop{\mathrm{lim}}\nolimits}
\def\colim{\mathop{\mathrm{colim}}\nolimits}
\def\Spec{\mathop{\mathrm{Spec}}}
\def\Hom{\mathop{\mathrm{Hom}}\nolimits}
\def\Ext{\mathop{\mathrm{Ext}}\nolimits}
\def\SheafHom{\mathop{\mathcal{H}\!\mathit{om}}\nolimits}
\def\SheafExt{\mathop{\mathcal{E}\!\mathit{xt}}\nolimits}
\def\Sch{\mathit{Sch}}
\def\Mor{\mathop{\mathrm{Mor}}\nolimits}
\def\Ob{\mathop{\mathrm{Ob}}\nolimits}
\def\Sh{\mathop{\mathit{Sh}}\nolimits}
\def\NL{\mathop{N\!L}\nolimits}
\def\CH{\mathop{\mathrm{CH}}\nolimits}
\def\proetale{{pro\text{-}\acute{e}tale}}
\def\etale{{\acute{e}tale}}
\def\QCoh{\mathit{QCoh}}
\def\Ker{\mathop{\mathrm{Ker}}}
\def\Im{\mathop{\mathrm{Im}}}
\def\Coker{\mathop{\mathrm{Coker}}}
\def\Coim{\mathop{\mathrm{Coim}}}

% Boxtimes
%
\DeclareMathSymbol{\boxtimes}{\mathbin}{AMSa}{"02}

%
% Macros for moduli stacks/spaces
%
\def\QCohstack{\mathcal{QC}\!\mathit{oh}}
\def\Cohstack{\mathcal{C}\!\mathit{oh}}
\def\Spacesstack{\mathcal{S}\!\mathit{paces}}
\def\Quotfunctor{\mathrm{Quot}}
\def\Hilbfunctor{\mathrm{Hilb}}
\def\Curvesstack{\mathcal{C}\!\mathit{urves}}
\def\Polarizedstack{\mathcal{P}\!\mathit{olarized}}
\def\Complexesstack{\mathcal{C}\!\mathit{omplexes}}
% \Pic is the operator that assigns to X its picard group, usage \Pic(X)
% \Picardstack_{X/B} denotes the Picard stack of X over B
% \Picardfunctor_{X/B} denotes the Picard functor of X over B
\def\Pic{\mathop{\mathrm{Pic}}\nolimits}
\def\Picardstack{\mathcal{P}\!\mathit{ic}}
\def\Picardfunctor{\mathrm{Pic}}
\def\Deformationcategory{\mathcal{D}\!\mathit{ef}}
\usepackage{cleveref}
\crefname{originallemma}{original lemma}{original lemmas}
\Crefname{originallemma}{Original lemma}{Original lemmas}
\makeatletter
\newlabel{lemma-local-homotopy}{{08FP}{source}{}{sourcepointer-1}{}}
\newlabel{lemma-local-homotopy@cref}{{[originallemma][0][]08FP}{[1][0][]source}{}{}{}}
\makeatother
\title{Lifting a map from a strictly perfect complex through a quasi-isomorphism, locally up to homotopy}
\author{The Stacks Project Authors}
\date{stacks-project:sites-cohomology.tex:local-snapshot-20261011; GNU FDL 1.2 or later}
\begin{document}
\maketitle
\noindent\textit{Editorial note.} All delivered mathematical text is the author's own, copied from the pinned Stacks Project chapter file sites-cohomology.tex; no statement or proof step was rewritten, reordered or removed.\par\medskip
\noindent\textit{Editorial note.} The proof is the author's complete argument, ending with the author's own closing phrase, which is retained verbatim.\par\medskip
\noindent\textit{Editorial note.} The Xy-pic diagram of the counted statement is delivered as the author wrote it, including the dotted arrow.\par\medskip
\noindent\textit{Editorial note.} Cross-chapter restatement: the same label occurs in chapter cohomology for a ringed space. That statement is the special case of this ringed-site statement, so per the lane-s4 rule only this more general occurrence is counted here.\par\medskip
\noindent\textit{Editorial note.} The single labelled reference occurring inside the counted spans is carried in assembly.external\_references as a non-load-bearing navigation pointer to the same chapter.\par\medskip
\noindent\textit{Editorial note.} The preamble is the chapter's own retained minimal preamble (sources/stacks-sites-cohomology/retained-preamble.tex), a verbatim excerpt of the native preamble: the Xy-pic setup and the native theorem declarations with the mathematical macros.\par\medskip
\noindent\textit{Allowed standard prerequisites (editorial).}\begin{itemize}
\item Strictly perfect complex of O\_U-modules. Same chapter, the definition of a strictly perfect complex as a bounded complex of direct summands of finite free modules. Supplies the hypothesis on E\textasciicircum{}bullet that makes the local triviality lemma applicable.
\item Cone of a morphism of complexes and its cohomology. Same chapter, the cone C(f)\textasciicircum{}bullet of a morphism f and the associated distinguished triangle. Supplies the reduction of the hypotheses on H\textasciicircum{}j(f) to the vanishing of the cohomology of the cone in degrees at least a.
\item Distinguished triangles in the homotopy category K(O\_U). Same chapter, the triangle G\textasciicircum{}bullet -> F\textasciicircum{}bullet -> C(f)\textasciicircum{}bullet -> G\textasciicircum{}bullet[1] and the long exact sequence of homotopy classes it induces. Supplies the final step, where vanishing of the composite is converted into a lift up to homotopy.
\end{itemize}\medskip
\ifdefined\hypertarget\hypertarget{sourcepointer-1}{}\fi
\noindent\textit{Source reference.} Original lemma 08FP: sources/stacks-sites-cohomology/pinned-original-source.tex, source lines 11796--11796. Same-chapter navigation pointer invoked at proof line 11866, where the author applies the local triviality of maps from a strictly perfect complex to a complex with vanishing cohomology above a. The pointer only resolves that reference for the reader; the cited lemma is not reproduced inside the counted span, and no bibliography citation is involved.\par\medskip
\begin{lemma}
\label{lemma-lift-through-quasi-isomorphism}
Let $(\mathcal{C}, \mathcal{O})$ be a ringed site. Let $U$ be an object of
$\mathcal{C}$. Given a solid diagram of complexes of $\mathcal{O}_U$-modules
$$
\xymatrix{
\mathcal{E}^\bullet \ar@{..>}[dr] \ar[r]_\alpha & \mathcal{F}^\bullet \\
& \mathcal{G}^\bullet \ar[u]_f
}
$$
with $\mathcal{E}^\bullet$ strictly perfect, $\mathcal{E}^j = 0$ for
$j < a$ and $H^j(f)$ an isomorphism for $j > a$ and surjective for
$j = a$, then there exists a covering $\{U_i \to U\}$ and for each $i$
a dotted arrow over $U_i$ making the diagram commute up to homotopy.
\end{lemma}

\begin{proof}
Our assumptions on $f$ imply the cone $C(f)^\bullet$ has vanishing
cohomology sheaves in degrees $\geq a$.
Hence Lemma \ref{lemma-local-homotopy} guarantees there is a
covering $\{U_i \to U\}$ such that the composition
$\mathcal{E}^\bullet \to \mathcal{F}^\bullet \to C(f)^\bullet$
is homotopic to zero over $U_i$. Since
$$
\mathcal{G}^\bullet \to \mathcal{F}^\bullet \to C(f)^\bullet \to
\mathcal{G}^\bullet[1]
$$
restricts to a distinguished triangle in $K(\mathcal{O}_{U_i})$
we see that we can lift $\alpha|_{U_i}$ up to homotopy to a map
$\alpha_i : \mathcal{E}^\bullet|_{U_i} \to \mathcal{G}^\bullet|_{U_i}$
as desired.
\end{proof}

\end{document}
